% =========================================================
% Bound Arithmetic
% Source: volume-iii/analysis/bounding/notes/notes-bounds-arithmetic.tex
% Slot: after notes-bounding-monotonicty.tex in notes/index.tex
% =========================================================

\subsection{Bound Arithmetic}
\label{sec:bound-arithmetic}

% ---------------------------------------------------------
% Toolkit
% ---------------------------------------------------------
\begin{tcolorbox}[colback=gray!6, colframe=gray!40, arc=2pt,
  left=6pt, right=6pt, top=4pt, bottom=4pt,
  title={\small\textbf{Bound Arithmetic — Quick Reference}},
  fonttitle=\small\bfseries]
\renewcommand{\arraystretch}{1.4}
\begin{tabular}{@{}p{0.22\textwidth} p{0.44\textwidth} p{0.26\textwidth}@{}}
\textbf{Result} & \textbf{Identity} & \textbf{Proof} \\
\midrule
Additivity (sup)
  & $\sup(A+B)=\sup A+\sup B$
  & \hyperref[prf:additivity-of-supremums]{Proof $\to$} \\
Additivity (inf)
  & $\inf(A+B)=\inf A+\inf B$
  & \hyperref[prf:additivity-of-infimums]{Proof $\to$} \\
Negation
  & $\sup(-A)=-\inf A,\quad\inf(-A)=-\sup A$
  & \hyperref[prf:supremum-infimum-swap]{Proof $\to$} \\
Scalar ($c>0$)
  & $\sup(cA)=c\sup A,\quad\inf(cA)=c\inf A$
  & \hyperref[prf:infimum-positive-scalar-multiple]{Proof $\to$} \\
Scalar ($c<0$)
  & $\sup(cA)=c\inf A,\quad\inf(cA)=c\sup A$
  & \hyperref[prf:supremum-infimum-swap]{Proof $\to$} \\
Difference (sup)
  & $\sup(A-B)=\sup A-\inf B$
  & \hyperref[prf:supremum-set-difference]{Proof $\to$} \\
Difference (inf)
  & $\inf(A-B)=\inf A-\sup B$
  & \hyperref[prf:supremum-set-difference]{Proof $\to$} \\
\end{tabular}
\end{tcolorbox}

\begin{remark*}
The definitions of supremum and infimum apply to a single fixed set.
When sets are combined — by addition, scaling, or subtraction — the
supremum of the resulting set must be recomputed from scratch unless a
reduction formula is available.
Bound arithmetic is that collection of formulas: six identities that let
bounds be manipulated the way numbers are.
They appear constantly in real analysis — in $\varepsilon$-$\delta$ arguments,
in integration, in convexity, in optimization.

The entire arithmetic descends from two load-bearing facts:
the \emph{additivity rule}
(Theorem~\ref{thm:sup-add}: $\sup(A+B) = \sup A + \sup B$)
and the \emph{negation rule}
(Theorem~\ref{thm:neg-rule}: $\sup(-A) = -\inf A$).
Every other identity reduces to applying these two in sequence.
\end{remark*}

% ---------------------------------------------------------
% Set arithmetic definitions
% ---------------------------------------------------------
\subsubsection{Set Arithmetic}

\begin{tcolorbox}[colback=propbox, colframe=propborder, arc=2pt,
  left=6pt, right=6pt, top=4pt, bottom=4pt,
  title={\small\textbf{Definition (Set Arithmetic)}},
  fonttitle=\small\bfseries]
\label{def:set-arithmetic}
Let $A, B \subseteq \mathbb{R}$ be nonempty and let $c \in \mathbb{R}$.
\begin{align*}
A + B &:= \{\, a + b \;:\; a \in A,\; b \in B \,\}
         && \text{(Minkowski sum)} \\
A - B &:= \{\, a - b \;:\; a \in A,\; b \in B \,\}
         && \text{(arithmetic difference)} \\
cA   &:= \{\, ca     \;:\; a \in A \,\}
         && \text{(scalar multiple)} \\
-A   &:= (-1)A = \{\, -a \;:\; a \in A \,\}
         && \text{(negation)}
\end{align*}
\end{tcolorbox}

\begin{remark*}[Fully quantified form]
\begin{align*}
x \in A+B  &\iff \exists\, a \in A\;\exists\, b \in B\;(x = a+b). \\
x \in A-B  &\iff \exists\, a \in A\;\exists\, b \in B\;(x = a-b). \\
x \in cA   &\iff \exists\, a \in A\;(x = ca). \\
x \in {-A} &\iff \exists\, a \in A\;(x = -a).
\end{align*}
\end{remark*}

\begin{remark*}[Notation]
The symbol $A - B$ is the \emph{arithmetic} difference — a new set whose
elements are all pairwise differences $a - b$ for $a \in A$, $b \in B$.
It is not the set-theoretic difference $\{x \in A : x \notin B\}$.
In this chapter $A - B$ always means the arithmetic version.
\end{remark*}

\begin{remark*}[$A - B = A + (-B)$]
\label{rem:diff-as-sum}
The fundamental rewriting rule is
\[
  A - B \;=\; A + (-B).
\]
This is not merely notation — it is the proof strategy for every
difference identity. Wherever $A - B$ appears, replace it with
$A + (-B)$, apply the additivity rule to the sum, then apply the negation
rule to $-B$. All difference identities cost no additional work.
\end{remark*}

\begin{remark*}[Concrete example]
Let $A = [1,3]$ and $B = [2,5]$.
\[
  A + B = [3,8], \qquad
  A - B = [-4,1], \qquad
  2A = [2,6], \qquad
  {-B} = [-5,-2].
\]
Reading off directly: $\sup(A+B) = 8 = 3 + 5 = \sup A + \sup B$.~$\checkmark$

The identities below explain why every such computation works
without enumerating the resulting set.
\end{remark*}

% ---------------------------------------------------------
% Master theorem
% ---------------------------------------------------------
\subsubsection{The Master Theorem}

\begin{tcolorbox}[colback=thmbox, colframe=thmborder, arc=2pt,
  left=6pt, right=6pt, top=4pt, bottom=4pt,
  title={\small\textbf{Theorem — Algebra of Suprema and Infima}},
  fonttitle=\small\bfseries]
\label{thm:bound-arithmetic-master}
Let $A, B \subset \mathbb{R}$ be nonempty, bounded above and below.

\medskip\noindent\textbf{Additivity.}
\begin{align}
\sup(A + B) &= \sup A + \sup B \label{eq:sup-add} \\
\inf(A + B) &= \inf A + \inf B \label{eq:inf-add}
\end{align}

\noindent\textbf{Scalar multiplication.} For $c \neq 0$:
\begin{align}
\sup(cA) &=
\begin{cases}
c\sup A & c > 0 \\[2pt]
c\inf A & c < 0
\end{cases}
\label{eq:sup-scale} \\[4pt]
\inf(cA) &=
\begin{cases}
c\inf A & c > 0 \\[2pt]
c\sup A & c < 0
\end{cases}
\label{eq:inf-scale}
\end{align}

\noindent\textbf{Differences.}
\begin{align}
\sup(A - B) &= \sup A - \inf B \label{eq:sup-diff} \\
\inf(A - B) &= \inf A - \sup B \label{eq:inf-diff}
\end{align}
\end{tcolorbox}

\begin{remark*}[Structure of the six identities]
The six identities split into three natural pairs: sums, scalar multiples,
differences. Within each pair the $\sup$ and $\inf$ statements mirror each
other exactly. That symmetry is the algebraic shadow of the order-reversal
caused by negation, examined carefully in the next section.
\end{remark*}

% ---------------------------------------------------------
% Geometric intuition
% ---------------------------------------------------------
\subsubsection{Geometric Intuition}

\paragraph{Negation reverses order.}

The simplest case: what does multiplying a set by $-1$ do to its bounds?

\begin{theorem}[\texorpdfstring{\hyperref[prf:supremum-infimum-swap]{Negation Rule}}{Negation Rule}]
\label{thm:neg-rule}
Let $A \subseteq \mathbb{R}$ be nonempty and bounded. Then
\[
\sup(-A) = -\inf A, \qquad \inf(-A) = -\sup A.
\]
\end{theorem}

\begin{remark*}[Fully quantified form]
$\forall\, A \subseteq \mathbb{R}$ nonempty and bounded:
$\sup\{-a : a \in A\} = -\inf A$
and
$\inf\{-a : a \in A\} = -\sup A$.
\end{remark*}

The picture is a reflection across $0$ on the number line.

\begin{figure}[h]
\centering
\begin{tikzpicture}[>=Stealth, thick, font=\small]
  \draw[->] (-6.5,0) -- (6.5,0) node[right] {$\mathbb{R}$};
  \draw (0,-0.18) -- (0,0.18) node[above]{\small $0$};

  % Set A = [1.5, 4.0]
  \draw[blue!70!black, line width=4pt, line cap=round] (1.5,0) -- (4.0,0);
  \node[blue!70!black, above] at (2.75,0.10) {$A$};
  \draw[blue!70!black, dashed, thin] (1.5,-0.38)--(1.5,0.68)
    node[above]{\small $\inf A$};
  \draw[blue!70!black, dashed, thin] (4.0,-0.38)--(4.0,0.68)
    node[above]{\small $\sup A$};

  % Set -A = [-4.0, -1.5]
  \draw[red!70!black, line width=4pt, line cap=round] (-4.0,0) -- (-1.5,0);
  \node[red!70!black, below] at (-2.75,-0.10) {$-A$};
  \draw[red!70!black, dashed, thin] (-4.0, 0.38)--(-4.0,-0.68)
    node[below]{\small $-\sup A = \inf(-A)$};
  \draw[red!70!black, dashed, thin] (-1.5, 0.38)--(-1.5,-0.68)
    node[below]{\small $-\inf A = \sup(-A)$};

  % Reflection arrows
  \draw[->, violet!70!black, bend right=30, thin]
    (4.0, 0.30) to node[above,font=\scriptsize]{$x \mapsto -x$} (-4.0,0.30);
  \draw[->, violet!70!black, bend right=30, thin]
    (1.5, 0.30) to (-1.5,0.30);
\end{tikzpicture}
\caption{Reflecting $A$ across $0$ maps $[\inf A,\sup A]$ to
$[-\sup A, -\inf A]$.
The right endpoint $\sup A$ maps to $-\sup A$, which is the
\emph{left} endpoint of $-A$ — hence $\inf(-A)$.
The left endpoint $\inf A$ maps to $-\inf A$, which is the
\emph{right} endpoint of $-A$ — hence $\sup(-A)$.}
\label{fig:negation-reflection}
\end{figure}

The figure makes the rule visceral. $A = [\inf A, \sup A]$ reflects to
$-A = [-\sup A, -\inf A]$: the interval flips entirely.
The right endpoint of $A$ — which was $\sup A$ — maps to $-\sup A$,
the \emph{left} endpoint of $-A$, and therefore becomes $\inf(-A)$.
The left endpoint of $A$ — which was $\inf A$ — maps to $-\inf A$,
the right endpoint of $-A$, and therefore becomes $\sup(-A)$.

\medskip\noindent
\textbf{Key take-away}: negation flips the sign \emph{and} swaps which
bound is being computed.
This is the engine behind every identity in this section.

\paragraph{The $2\times 2$ symmetry diagram.}

The negation rule compresses into the following diagram:

\begin{figure}[h]
\centering
\begin{tikzpicture}[>=Stealth, node distance=1.8cm and 4.2cm,
    every node/.style={font=\small}]
  \node (supA)  {$\sup A$};
  \node (infA)  [below of=supA] {$\inf A$};
  \node (ninfA) [right of=supA] {$-\inf A = \sup(-A)$};
  \node (nsupA) [right of=infA] {$-\sup A = \inf(-A)$};

  \draw[->] (supA)  -- node[above]{\small $\times(-1)$} (ninfA);
  \draw[->] (infA)  -- node[below]{\small $\times(-1)$} (nsupA);
  \draw[<->] (supA) -- node[left]{\small swap} (infA);
  \draw[<->] (ninfA)-- node[right]{\small swap} (nsupA);

  \node[above=0.15cm of supA,  font=\footnotesize, gray] {on $A$};
  \node[above=0.15cm of ninfA, font=\footnotesize, gray] {on $-A$};
\end{tikzpicture}
\caption{Horizontal arrows: multiply by $-1$.
Vertical arrows: swap $\sup \leftrightarrow \inf$.
The diagram commutes: go right then down, or go down then right —
the same corner is reached either way.}
\label{fig:negation-symmetry}
\end{figure}

The horizontal arrows represent multiplication by $-1$.
The vertical movement represents switching between upper and lower bounds.
Everything is consistent: go right, then down, or go down, then right —
the same corner is reached either way.
A negated bound cannot be written down without also swapping which bound
is being computed.
This diagram should remain visible when working through any of the proofs
that follow.

\paragraph{Geometry of the sum and difference.}

The figure below illustrates the proof strategy for
$\sup(A-B) = \sup A - \inf B$.
Each row is a step in the argument.

\begin{figure}[h]
\centering
%% Colour palette (matches project house style: muted, academic)
\definecolor{colA}{RGB}{70,130,180}    % steel blue  — set A
\definecolor{colB}{RGB}{200,80,60}     % brick red   — set B
\definecolor{colC}{RGB}{80,160,100}    % sage green  — C = -B
\definecolor{colD}{RGB}{130,80,180}    % purple      — A - B = A + C

 
\begin{tikzpicture}[
  >=Stealth,
  thick,
  font=\small,
  every label/.style={font=\footnotesize},
]
 
%% ---------------------------------------------------------------
%% Layout constants
%% ---------------------------------------------------------------
\def\rowA{0}      % y-coord of Row 1: sets A and B
\def\rowB{-2.4}   % y-coord of Row 2: C = -B
\def\rowC{-4.8}   % y-coord of Row 3: A - B = A + C
 
%% ---------------------------------------------------------------
%% ROW 1 — Sets A and B on the real line
%% ---------------------------------------------------------------
 
%% ---- Axis for row 1 ----
\draw[->] (-0.5,\rowA) -- (10.5,\rowA) node[right] {$\mathbb{R}$};
 
%% ---- Set A: interval [1.5, 4.0] ----
\def\Alo{1.5}  \def\Ahi{4.0}
\def\supA{4.0}
\draw[colA, line width=4pt, line cap=round]
  (\Alo,\rowA) -- (\Ahi,\rowA);
\node[colA, above] at ({(\Alo+\Ahi)/2},\rowA+0.05) {$A$};
% sup A marker
\draw[colA, dashed, thin] (\supA,\rowA-0.25) -- (\supA,\rowA+0.55);
\node[colA, above] at (\supA,\rowA+0.55) {$\sup A$};
 
%% ---- Set B: interval [5.5, 8.0] ----
\def\Blo{5.5}  \def\Bhi{8.0}
\def\infB{5.5}
\draw[colB, line width=4pt, line cap=round]
  (\Blo,\rowA) -- (\Bhi,\rowA);
\node[colB, above] at ({(\Blo+\Bhi)/2},\rowA+0.05) {$B$};
% inf B marker
\draw[colB, dashed, thin] (\infB,\rowA-0.25) -- (\infB,\rowA+0.55);
\node[colB, above] at (\infB,\rowA+0.55) {$\inf B$};
 
%% Row 1 label
\node[left, align=right] at (-0.5,\rowA) {\textbf{Step 0}\\$A,B \subset \mathbb{R}$};
 
%% ---------------------------------------------------------------
%% ROW 2 — C = -B (reflection of B through 0)
%%          We shift the axis so 0 lines up nicely.
%%          B lives at [5.5, 8.0]; -B lives at [-8.0, -5.5].
%%          We offset the display: map coordinate x -> x + 9 for row 2.
%% ---------------------------------------------------------------
\def\shift{9}   % display offset so -B appears in visible range
 
\draw[->] (-0.5,\rowB) -- (10.5,\rowB) node[right] {$\mathbb{R}$};
 
%% C = -B drawn at [9-8, 9-5.5] = [1, 3.5]
\def\Clo{1.0}  \def\Chi{3.5}
\def\supC{3.5}
\draw[colC, line width=4pt, line cap=round]
  (\Clo,\rowB) -- (\Chi,\rowB);
\node[colC, above] at ({(\Clo+\Chi)/2},\rowB+0.05) {$C = {-B}$};
% sup C = -inf B
\draw[colC, dashed, thin] (\supC,\rowB-0.25) -- (\supC,\rowB+0.55);
\node[colC, above] at (\supC,\rowB+0.55) {$\sup C = {-\inf B}$};
 
%% Reflection arrow annotation
\draw[->, colB!60!colC, thin, bend left=25]
  (\Blo-0.1, \rowA-0.18)
  to node[midway, right, font=\scriptsize, colB!60!colC]
     {$b \mapsto -b$}
  (\supC+0.1, \rowB+0.18);
 
\node[left, align=right] at (-0.5,\rowB)
  {\textbf{Step 1}\\$C = {-B}$};
 
%% ---------------------------------------------------------------
%% ROW 3 — A + C  (Minkowski sum: sup = sup A + sup C)
%%
%%  sup(A+C) = sup A + sup C = 4.0 + 3.5 = 7.5
%%  The set A+C = [Alo+Clo, Ahi+Chi] = [2.5, 7.5]
%%  Scale to fit: divide everything by ~1 (they fit already)
%% ---------------------------------------------------------------
\draw[->] (-0.5,\rowC) -- (10.5,\rowC) node[right] {$\mathbb{R}$};
 
%% A + C = [1.5+1.0, 4.0+3.5] = [2.5, 7.5]
\def\Dlo{2.5}  \def\Dhi{7.5}
\def\supD{7.5}
\draw[colD, line width=4pt, line cap=round]
  (\Dlo,\rowC) -- (\Dhi,\rowC);
\node[colD, above] at ({(\Dlo+\Dhi)/2},\rowC+0.05) {$A - B = A + C$};
\draw[colD, dashed, thin] (\supD,\rowC-0.25) -- (\supD,\rowC+0.55);
\node[colD, above] at (\supD,\rowC+0.55)
  {$\sup(A{-}B) = \sup A - \inf B$};
 
%% ---- Additive decomposition brace ----
%% Show: the distance from Dlo to Dhi splits into (supA - Alo) and (Ahi - something)
%% Better: show sup(A+C) = supA + supC with a small split bar at supA (=4.0)
%% Mark sup A contribution and sup C contribution along row 3
\def\splitpt{4.0}   % sup A value, and also A's contribution starts at Dlo
%% Brace under [Dlo, splitpt]: "from A" part
\draw[decorate, decoration={brace, amplitude=5pt, mirror},
      colA, thin]
  (\Dlo, \rowC-0.35) -- (\splitpt, \rowC-0.35)
  node[midway, below=6pt, colA, font=\scriptsize]
    {spread of $A$};
 
%% Brace under [splitpt, Dhi]: "from C" part
\draw[decorate, decoration={brace, amplitude=5pt, mirror},
      colC, thin]
  (\splitpt, \rowC-0.35) -- (\Dhi, \rowC-0.35)
  node[midway, below=6pt, colC, font=\scriptsize]
    {spread of $C$};
 
\node[left, align=right] at (-0.5,\rowC)
  {\textbf{Step 2}\\$A{+}C$};
 
%% ---------------------------------------------------------------
%% Vertical alignment guides: drop faint lines from sup A and inf B
%% down through all rows to show how the values propagate.
%% ---------------------------------------------------------------
\draw[colA!40, very thin, dashed]
  (\supA, \rowA-0.3) -- (\supA, \rowC+0.15);
\draw[colB!40, very thin, dashed]
  (\infB, \rowA-0.3) -- (\infB, \rowC+0.15);   % infB=5.5 passes through all rows
 
%% ---------------------------------------------------------------
%% Title
%% ---------------------------------------------------------------
\node[font=\normalsize\bfseries] at (5, 1.3)
  {Geometry of $\sup(A - B) = \sup A - \inf B$};
 
\end{tikzpicture}
\caption{%
\textbf{Step~0.} The original sets $A$ (blue) and $B$ (red) on $\mathbb{R}$.
Dashed vertical guides track $\sup A$ and $\inf B$ downward through all rows.
\textbf{Step~1.} Form $C = -B$ (green). Reflection sends the left endpoint
of $B$ (which is $\inf B$) to the right endpoint of $C$,
so $\sup C = -\inf B$. The vertical guide confirms the alignment.
\textbf{Step~2.} The Minkowski sum $A + C$ (purple) has
$\sup(A+C) = \sup A + \sup C = \sup A - \inf B$.
The underbrace decomposes the right endpoint: it splits into the
contribution from $A$ and the contribution from $C$.}
\label{fig:sup-diff-geometry}
\end{figure}

% ---------------------------------------------------------
% Additivity
% ---------------------------------------------------------
\subsubsection{Additivity}

\begin{theorem}[\texorpdfstring{\hyperref[prf:additivity-of-supremums]{Additivity of Supremum}}{Additivity of Supremum}]
\label{thm:sup-add}
Let $A, B \subseteq \mathbb{R}$ be nonempty and bounded above. Then
\[
\sup(A+B) = \sup A + \sup B.
\]
\end{theorem}

\begin{theorem}[\texorpdfstring{\hyperref[prf:additivity-of-infimums]{Additivity of Infimum}}{Additivity of Infimum}]
\label{thm:inf-add}
Let $A, B \subseteq \mathbb{R}$ be nonempty and bounded below. Then
\[
\inf(A+B) = \inf A + \inf B.
\]
\end{theorem}

\begin{remark*}[Fully quantified form]
$\forall\, A, B \subseteq \mathbb{R}$ nonempty, bounded above:
$\sup\{a + b : a \in A,\, b \in B\} = \sup A + \sup B$.
\end{remark*}

\begin{remark*}[Why the proof needs two parts]
Every upper-bound proof requires two parts.
\textbf{Part~1} shows the candidate is an upper bound.
\textbf{Part~2} shows it is the \emph{least} one — that nothing smaller
works. Omitting Part~2 establishes only
$\sup(A+B) \le \sup A + \sup B$, which is strictly weaker than equality.

For additivity: Part~1 is a single line
($a + b \le \sup A + \sup B$ follows immediately from adding the two
defining inequalities).
Part~2 uses the $\varepsilon$-characterization
(Proposition~\ref{prop:eps-char}):
set $\varepsilon = (s + t - u)/2 > 0$ for any candidate $u < s + t$,
extract $a \in A$ with $a > \sup A - \varepsilon$ and
$b \in B$ with $b > \sup B - \varepsilon$, and add to get
$a + b > s + t - 2\varepsilon = u$.
So $u$ is not an upper bound, and $s + t = \sup(A+B)$.

The infimum version is the identical argument with all inequalities reversed:
wherever $a \le s$ appeared, use $a \ge \inf A$;
where an element exceeding $u$ was found, find one below $u$.
The logic is a mirror image.
\end{remark*}

% ---------------------------------------------------------
% Negation rule (proof structure remark belongs here)
% ---------------------------------------------------------
\begin{remark*}[Negation rule: why it is the key to everything]
Once the negation rule is in hand, the scalar and difference identities
cost almost nothing.
Every one of them reduces to applying
Theorem~\ref{thm:neg-rule} once and invoking additivity.
The proof of the negation rule itself is a two-part upper-bound argument
on $-A$: show $-\inf A$ is an upper bound, then show it is the least one
by pushing any smaller candidate below $\inf A$.
The second identity $\inf(-A) = -\sup A$ follows by applying the first
to $-A$:
$\sup(A) = \sup(-(-A)) = -\inf(-A)$, so $\inf(-A) = -\sup A$.
\end{remark*}

% ---------------------------------------------------------
% Scalar multiplication
% ---------------------------------------------------------
\subsubsection{Scalar Multiplication}

\begin{theorem}[\texorpdfstring{\hyperref[prf:infimum-positive-scalar-multiple]{Scalar Multiplication ($c > 0$)}}{Scalar Multiplication (c > 0)}]
\label{thm:scale-pos}
Let $A \subseteq \mathbb{R}$ be nonempty and bounded, and let $c > 0$. Then
\[
\sup(cA) = c\sup A, \qquad \inf(cA) = c\inf A.
\]
\end{theorem}

\begin{theorem}[\texorpdfstring{\hyperref[prf:supremum-infimum-swap]{Scalar Multiplication ($c < 0$)}}{Scalar Multiplication (c < 0)}]
\label{thm:scale-neg}
Let $A \subseteq \mathbb{R}$ be nonempty and bounded, and let $c < 0$. Then
\[
\sup(cA) = c\inf A, \qquad \inf(cA) = c\sup A.
\]
\end{theorem}

\begin{remark*}[Fully quantified form]
$\forall\, c > 0$, $\forall\, A \subseteq \mathbb{R}$ nonempty, bounded above:
$\sup\{ca : a \in A\} = c\sup A$.
$\forall\, c < 0$, $\forall\, A \subseteq \mathbb{R}$ nonempty, bounded above:
$\sup\{ca : a \in A\} = c\inf A$.
\end{remark*}

\begin{remark*}[Why the cases split]
Multiplying by $c > 0$ stretches the set without reflection: the ordering
of all elements is preserved, so the right endpoint stays the right
endpoint and simply gets scaled.
Multiplying by $c < 0$ reflects the set across $0$ and then stretches:
the right endpoint of $A$ becomes the left endpoint of $cA$.
So $\sup A$ maps to $c\sup A = \inf(cA)$ and $\inf A$ maps to
$c\inf A = \sup(cA)$.

For the proof: the $c < 0$ case is a one-line reduction.
Write $c = -|c|$, so $cA = -(|c|A)$.
Apply the $c > 0$ result to $|c|A$: $\sup(|c|A) = |c|\sup A$.
Then apply the negation rule:
$\sup(cA) = \sup(-(|c|A)) = -\inf(|c|A) = -|c|\inf A = c\inf A$.
No new ideas are needed — only the two master facts chained.
\end{remark*}

\begin{figure}[h]
\centering
\begin{tikzpicture}[>=Stealth, thick, font=\small]

  %% Row 1: A (blue) and 2A (green)
  \draw[->] (-0.3,1.6) -- (10.2,1.6);
  \draw[blue!70!black, line width=4pt, line cap=round] (1,1.6) -- (3,1.6);
  \node[blue!70!black, above] at (2,1.68) {$A$};
  \draw[blue!70!black, dashed, thin] (1,1.3)--(1,2.35)
    node[above]{\scriptsize $\inf A$};
  \draw[blue!70!black, dashed, thin] (3,1.3)--(3,2.35)
    node[above]{\scriptsize $\sup A$};
  \draw[green!50!black, line width=4pt, line cap=round] (2,1.6) -- (6,1.6);
  \node[green!50!black, above] at (4.2,1.68) {$2A$};
  \draw[green!50!black, dashed, thin] (6,1.3)--(6,2.35)
    node[above]{\scriptsize $2\sup A = \sup(2A)$};
  \node[right, font=\footnotesize] at (10.2,1.6) {$c=2>0$: stretches};

  %% Row 2: A (blue) and -2A (red)
  \draw[->] (-7.5,0) -- (4.8,0);
  \draw (0,-0.18)--(0,0.18) node[above]{\scriptsize $0$};
  \draw[blue!70!black, line width=4pt, line cap=round] (1,0) -- (3,0);
  \node[blue!70!black, above] at (2,0.08) {$A$};
  \draw[blue!70!black, dashed, thin] (1,-0.42)--(1,0.58)
    node[above]{\scriptsize $\inf A$};
  \draw[blue!70!black, dashed, thin] (3,-0.42)--(3,0.58)
    node[above]{\scriptsize $\sup A$};
  \draw[red!70!black, line width=4pt, line cap=round] (-6,0) -- (-2,0);
  \node[red!70!black, below] at (-4,-0.08) {$-2A$};
  \draw[red!70!black, dashed, thin] (-2,-0.42)--(-2,0.58)
    node[above]{\scriptsize $-2\inf A = \sup(-2A)$};
  \draw[red!70!black, dashed, thin] (-6,-0.42)--(-6,0.58)
    node[above]{\scriptsize $-2\sup A = \inf(-2A)$};
  \draw[->, violet!70!black, bend right=28, thin]
    (3,0.38) to node[above,font=\scriptsize]{$\times(-2)$} (-6,0.38);
  \draw[->, violet!70!black, bend right=28, thin]
    (1,0.38) to (-2,0.38);
  \node[right, font=\footnotesize] at (4.8,0) {$c=-2<0$: reflects and stretches};

\end{tikzpicture}
\caption{Top: $c = 2 > 0$ stretches rightward; $\sup(2A) = 2\sup A$.
Bottom: $c = -2 < 0$ reflects across $0$ and stretches;
the old right endpoint $\sup A$ maps to the new left endpoint
$-2\sup A = \inf(-2A)$,
while the old left endpoint $\inf A$ maps to the new right endpoint
$-2\inf A = \sup(-2A) = c\inf A$.}
\label{fig:scalar-multiple}
\end{figure}

\begin{remark*}[Common error]
For $c < 0$: $\sup(cA) = c\,\inf A$, not $c\,\sup A$.
The sign change always comes with a bound-swap.
Counterexample: $A = [1,2]$, $c = -1$.
Then $-A = [-2,-1]$, so $\sup(-A) = -1 = (-1)\cdot\inf A$.
The formula $c\sup A = (-1)(2) = -2$ gives the \emph{infimum}, not the supremum.
\end{remark*}

% ---------------------------------------------------------
% Set difference
% ---------------------------------------------------------
\subsubsection{Set Differences}

\begin{theorem}[\texorpdfstring{\hyperref[prf:supremum-set-difference]{Supremum of a Difference}}{Supremum of a Difference}]
\label{thm:sup-diff}
Let $A, B \subseteq \mathbb{R}$ be nonempty, bounded above and below. Then
\[
\sup(A-B) = \sup A - \inf B, \qquad \inf(A-B) = \inf A - \sup B.
\]
\end{theorem}

\begin{remark*}[Fully quantified form]
$\forall\, A, B \subseteq \mathbb{R}$ nonempty and bounded:
$\sup\{a - b : a \in A,\, b \in B\} = \sup A - \inf B$.
\end{remark*}

\begin{remark*}[Intuition]
To maximise $a - b$: make $a$ as large as possible and $b$ as small as
possible.
So $\sup(A-B)$ is approached in the limit as $a \to \sup A$ and
$b \to \inf B$, giving $\sup A - \inf B$.
The formula $\sup(A-B) = \sup A - \sup B$ is the standard error:
it conflates maximising a difference with subtracting two suprema.
The geometry of Figure~\ref{fig:sup-diff-geometry} shows exactly why the
correct formula involves $\inf B$, not $\sup B$.
\end{remark*}

\begin{remark*}[The derivation in one chain]
The proof requires no new ideas. Every step applies a result already in hand:
\[
\sup(A-B)
\;=\; \sup\bigl(A+(-B)\bigr)
\;=\; \sup A + \sup(-B)
\;=\; \sup A + (-\inf B)
\;=\; \sup A - \inf B.
\]
Four equalities, three tools: rewrite as Minkowski sum
(Remark~\ref{rem:diff-as-sum}),
apply additivity \eqref{eq:sup-add},
apply the negation rule (Theorem~\ref{thm:neg-rule}).
If any step feels unearned, return to Figure~\ref{fig:sup-diff-geometry}:
each equality corresponds to a visible transformation on the number line.

For the infimum:
\[
\inf(A-B) = \inf A + \inf(-B) = \inf A + (-\sup B) = \inf A - \sup B.
\]
\end{remark*}

% ---------------------------------------------------------
% Unified summary
% ---------------------------------------------------------
\subsubsection{Unified Summary}

\begin{remark*}[The dependency tree]
Every identity ultimately rests on two roots:
\begin{enumerate}
  \item \textbf{Additivity}
        (Theorem~\ref{thm:sup-add}):
        $\sup(A+B) = \sup A + \sup B$.
  \item \textbf{Negation rule}
        (Theorem~\ref{thm:neg-rule}):
        $\sup(-A) = -\inf A$.
\end{enumerate}
The negation rule is completeness-free — it uses only order-reversal under
$\times(-1)$, which holds in any ordered set.
The additivity rule uses completeness indirectly, through the
$\varepsilon$-characterization (Proposition~\ref{prop:eps-char}).
All scaling ($c < 0$) and difference identities reduce to chaining these two.
\end{remark*}

\begin{tcolorbox}[colback=gray!6, colframe=gray!40, arc=2pt,
  left=6pt, right=6pt, top=4pt, bottom=4pt,
  title={\small\textbf{All identities and dependencies}},
  fonttitle=\small\bfseries]
\renewcommand{\arraystretch}{1.5}
\small
\begin{tabular}{@{}p{0.38\textwidth} p{0.22\textwidth} p{0.32\textwidth}@{}}
\textbf{Identity} & \textbf{Condition} & \textbf{Depends on} \\
\midrule
$\sup(A+B)=\sup A+\sup B$  & $A,B$ bounded above & direct \\
$\inf(A+B)=\inf A+\inf B$  & $A,B$ bounded below & direct (mirror) \\
$\sup(-A)=-\inf A$         & $A$ bounded         & direct \\
$\inf(-A)=-\sup A$         & $A$ bounded         & negation (sup) \\
$\sup(cA)=c\sup A$         & $c>0$               & direct scaling \\
$\inf(cA)=c\inf A$         & $c>0$               & direct scaling \\
$\sup(cA)=c\inf A$         & $c<0$               & negation $+$ scaling \\
$\inf(cA)=c\sup A$         & $c<0$               & negation $+$ scaling \\
$\sup(A-B)=\sup A-\inf B$  & $A,B$ bounded       & additivity $+$ negation \\
$\inf(A-B)=\inf A-\sup B$  & $A,B$ bounded       & additivity $+$ negation \\
\end{tabular}
\end{tcolorbox}

\begin{remark*}[Common errors]
\textbf{Wrong formula for $\sup(A-B)$.}
The temptation is $\sup A - \sup B$ by analogy with scalar arithmetic.
The correct formula is $\sup A - \inf B$.
To maximise $a - b$, want $a$ as large as possible and $b$ as small as
possible — hence $\sup A$ paired with $\inf B$, not $\sup B$.

\medskip\noindent
\textbf{Forgetting the bound-swap under negation.}
Negating a set maps $\sup \to -\inf$ and $\inf \to -\sup$.
Writing $\sup(-A) = -\sup A$ drops the swap.
Figure~\ref{fig:negation-reflection} shows exactly why: the right endpoint
of $A$ maps to the \emph{left} endpoint of $-A$.

\medskip\noindent
\textbf{Wrong bound in the negative-scalar rule.}
For $c < 0$: $\sup(cA) = c\,\inf A$, not $c\,\sup A$.
The sign change always comes with a bound-swap, for the same reason as
negation.
All three errors are variants of the same mistake: forgetting that
negation swaps which bound is being computed.
\end{remark*}

\begin{remark*}[Connections forward]
These identities appear throughout the rest of the notes:
\begin{itemize}
  \item \emph{Convergence}: bounding $|a_n + b_n|$ reduces to
        $\sup|a_n| + \sup|b_n|$ via additivity.
  \item \emph{Darboux integration}: $U(f+g,P) \le U(f,P) + U(g,P)$
        follows from additivity applied pointwise to the image of $f+g$.
  \item \emph{Oscillation}: $\operatorname{osc}(f,I) = \sup_I f - \inf_I f$
        is the difference identity applied to the image of $f$ on $I$.
  \item \emph{Metric spaces}: diameter estimates on $A - B$ use
        $\sup(A-B) = \sup A - \inf B$; see the monotonicity corollaries
        in Section~\ref{sec:bound-monotone}.
\end{itemize}

Once these patterns are fluent, many estimates in real analysis —
bounding $|f(x) - g(x)|$, controlling the $\sup$ of a sequence of
functions, estimating Darboux sums — reduce to one of these six identities.
\end{remark*}